% 本文件由 LaTeX 排版 Agent 根据有效 translation.json 自动生成。
% author=John Chrysostom; work_id=1908; title=给慕道者的训导

\chapter{给慕道者的训导}

% structure: p0001 source=h1 target=omitted reason=page-title-already-rendered-by-parent

% structure: p0002 source=h2 target=section reason=relative-heading-rank; base=h2

\section{第一训导}

致即将受光照者；并说明为何称洗礼池为重生之洗而非罪赦之洗；以及不仅背誓，即便真诚发誓也是危险的事。\par

1. 我们的年轻弟兄们的团体是多么令人喜悦和可爱！因为我称你们为弟兄，甚至是在你们尚未被产生之前，在你们出生之前，我就接纳这份与你们的关系：因为我知道，我清楚地知道，你们将如何被引向何等大的尊荣，何等大的尊严；而那些即将领受尊严的人，所有人通常都会在尊严授予之前就尊敬他们，藉着他们的关注为自己预先积存未来的善意。我现在也正是这样做。因为你们不是要被引向一个空虚的尊严，而是真正的王国；而且不仅仅是王国，而是天国本身。因此，我恳求并请求你们，当你们进入那个王国时，请记得我，如同若瑟对司酒长所说的：\Quote{当您得福时，请记得我。} \BibleRef{创 40:14} 现在我也对你们说这话：当你们得福时，请记得我。我要求这个，并非像他那样为报答你们解梦；因为我并非来为你们解梦，而是来谈论天上的事，并给你们带来如此美善的好消息，正如\Quote{眼所未见，耳所未闻，人心所未想到的，就是天主为爱祂的人所准备的。} \BibleRef{格前 2:9} 的确，若瑟对那司酒长说：\Quote{三天以内，法郎要恢复你的司酒长职位。} 但我说的不是三天以内你们将被委任去倒暴君的酒，而是三十天以后，不是法郎，而是天主要恢复你们在高天之上的家乡——自由的耶路撒冷——那在天上的城邑；而\Emph{他}确实说过：\Quote{你要把杯递到法郎手中。} 但我不是说你们要将杯递到王手中，而是王要将杯递到你们手中——那令人敬畏的杯，充满大能力，比任何受造物更珍贵。已经受洗者知道这杯的德能，你们不久也将知道。所以，当你们进入那个王国时，当你们穿上王家礼服时，当你们束上浸染着主血的紫袍时，当你们被戴上冠冕，其光芒四射，比太阳的光线更加灿烂时，请记得我。这些都是新郎的恩赐，确实大于你们的功德，却与祂的慈爱相称。\par

因此，我视你们在那些神圣婚筵之前便已蒙福，而且我不仅视你们为有福，还称赞你们的明智，因为你们不像那些最懒惰的人，在最后一口气才来接受光照，而是像明智的仆人，怀着极大的善意准备服从主人，以极大的温顺和甘心将你们灵魂的颈项置于基督的轭下，接受了祂轻松的轭和祂轻省的担子。因为如果所赐的恩宠对你们和那些临终时才受洗的人是一样的，然而意向的准备和对礼仪的准备却不相同。因为他们是在病床上领受，而你们是在教会——我们众人共同的母亲——的怀抱中领受；他们是在哀叹和哭泣中，而你们是在喜乐和极大欢欣中；他们叹息，你们感恩；他们因高热而昏沉，你们却充满属灵的喜乐。因此，在你们的情形中，一切与恩赐协调一致；在他们那里，一切却与恩赐相悖。因为受洗者那里有哀号和极大的悲叹，孩子们围立哭泣，妻子撕扯脸颊，朋友沮丧，仆人流涕；整个家庭景象如同一个阴沉冬日。如果你打开躺在那里的病人的心，你会发现它更忧郁。因为如同风以许多相反的冲击相遇，将海搅得四分五裂，同样，那吞噬他的恐惧思想也袭击病人的灵魂，以诸多焦虑分散其心神。每当他看见自己的孩子，就想到他们将成孤儿；每当他从他们转向妻子，就想着她的寡居；每当他看见仆人们，就看到整个家庭的荒凉；每当他回到自己，就想到自己当前的生命，即将与之分离，便经历到极大的沮丧阴云。那即将受洗者的灵魂就是如此。然后在它的骚动和混乱中，司铎进来了，他比热病本身更可怕，比死亡更令病人的亲属痛苦。因为司铎的进入被认为是比医生对生命绝望的声音更大的绝望理由，而那暗示永生的事物似乎成了死亡的象征。但我尚未给这些不幸画上句号。因为在亲属骚动和准备的喧嚣中，灵魂常常已经飞逝，留下孤零零的尸体；在许多情况下，即使灵魂还在，也是无用的，因为当它既不认识在场的人，也听不见他们的声音，不能回答那些使它与我们众人共同的主人订立那有福盟约的话语，而是如同一根无用的木头或一块石头，那即将受光照的人躺着与尸体无异，在这种情况下，受洗的益处何在？\par

2. 凡将接近这些神圣可畏之奥秘者，必须清醒警觉，免除今生一切挂虑，充满极大的节制与准备；他须摒除心中任何与奥秘相违之念，从各方面洁净并预备自己的居所，如同即将迎见君王本人一般。你心之准备如此；你思念如此；你灵魂之志向亦如此。故此，期望从天主获得与这最卓越决意相称的回报——祂以丰厚的酬报压倒那些对祂表现顺服的人。然而，由于祂的同仆也必须贡献自己的一份，那么我们也当贡献自己的一份；其实，连这些也不是我们自己的，而是我们主人的。\Quote{因为你有什么，}祂说，\Quote{不是领受的呢？既然是领受的，为什么你夸耀，好像不是领受的呢？} \BibleRef{格前 4:7} 我首先愿说明：为何我们的教父们，跳过整年，定于这个时节使教会子女领受入门圣事；以及为何在领受我们的教导之后，除去你们的鞋履和衣服，赤脚裸身，只着一件外衣，带领你们去听驱魔者的言语。因为他们并非轻率或鲁莽地为我们安排这衣饰和这个时节，而是这两者都有某种神秘与隐秘的理由。我愿将此事告诉你们。但我看到我们现在的讲论迫使我们必须谈及更必要的事，即什么是圣洗，它为何进入我们的生命，以及它给我们带来何等恩惠。\par

但若你愿意，让我们讨论这神秘净化的名称：它的名称并非单一，而是众多且各异的。因为这种净化被称为\Quote{重生之洗}。\Quote{祂拯救了我们，}他说，\Quote{借着重生之洗和圣神的更新。} \BibleRef{铎 3:5} 它也被称为\Quote{光照}，圣保禄也曾这样称呼：\Quote{请你们回想先前的日子，那时你们受了光照，忍受了极大的苦难斗争；} \BibleRef{希 10:32} 又说：\Quote{因为那些曾一次蒙受光照，尝过天上的恩赐，而后又离弃了的人，是不可能再使他们恢复悔改的。} \BibleRef{希 6:4} 它也被称为\Quote{圣洗}：\Quote{因为你们凡受过洗归于基督的，就是穿上了基督。} \BibleRef{迦 3:27} 它也被称为\Quote{埋葬}：\Quote{我们被埋葬了，}他说，\Quote{与他一起，借着圣洗，归于死亡。} \BibleRef{罗 6:4} 它也被称为\Quote{割损}：\Quote{你们也在祂内受了不是人手所行的割损，而是基督的割损，在于脱去肉欲的身体。} \BibleRef{哥 2:11} 它也被称为\Quote{十字架}：\Quote{我们的旧人已与祂同钉十字架，使罪恶的身体失效。} \BibleRef{罗 6:6} 此外还可以列举其他名称。但为避免我们将全部时间花费在这项恩赐的名称上，来吧，让我们回到第一个名称，并以解释其意义来结束讲论；同时，让我们将教导再稍加延伸。有众人共用的沐浴之洗，常洗去身体的不洁；又有犹太人的洗，比前者更尊贵，却远逊于恩宠之洗；它也洗去身体的不洁，但并非单纯的身体不洁，因为它甚至触及软弱的良心。因为有许多事物，按其本性并非不洁，但因良心的软弱而成为不洁。正如小儿害怕面具和其他鬼怪，这些事物本身并不可怕，但因小儿天性的软弱而显得可怕；我所论及之事也是如此：触摸尸体按其本性并非不洁，但若接触软弱的良心，便使触摸者成为不洁。因为此事按其本性并非不洁，连订立此法本身的梅瑟也证明了：他扛起若瑟的全尸，却仍保持洁净。为此，保禄论及这种并非来自本性而是因良心软弱而成的不洁时，这样说：\Quote{凡物本身没有不洁的，除非有人以为不洁，那物对他来说便是不洁的。} \BibleRef{罗 14:14} 你不见不洁并非出于事物本性，而是出于对这事判断的软弱吗？又说：\Quote{一切固然是洁净的，但若有人带着反感而吃，那对他就是恶了。} \BibleRef{罗 14:20} 你不见并非吃本身，而是带着反感吃，才是不洁的原因吗？\par

3. 犹太人的洗池所洁净的就是这样的污秽。但恩宠的洗池却并非如此，它洁净的是那将污秽引入灵魂和身体的真正不洁。因为它洁净的不是触摸死尸的人，而是着手于死行的人：若有任何人淫秽、行淫、拜偶像，或行任何你所能想到的恶事，或充满人间一切邪恶，一旦他落入这水洗池中，他从这神圣泉源再上来时，比太阳的光芒还要纯洁。为使你们不认为所说的只是空夸，请听\Person{保禄}论及这洗池的能力：\BibleQuote{你们不要自欺：无论是淫荡的、拜偶像的、奸淫的、作娈童的、行淫的、贪婪的、酗酒的、辱骂人的、勒索人的，都不能承继天主的国。}\BibleRef{格前 6:9}有人会说：\Quote{这与所说的话有什么相干？因为这证明洗池的能力是否彻底洗净这一切。}请听下文：\BibleQuote{你们中从前也有人是这样，但你们已经洗净了，已经圣化了，已经因主耶稣基督的名和我们天主的圣神成义了。}\Emph{我们}应许向你们显示：凡接近这洗池的人，都从一切淫乱中变得洁净；但这话更显示出：他们不仅洁净，而且既成圣又成义；因为它不只说\BibleQuote{你们洗净了}，还说\BibleQuote{你们圣化了，成义了}。还有什么比这更不可思议呢？不须劳苦、不须奋斗、不须善行，而正义油然而生。因为神圣恩赐的慈爱就是如此，它使人不需这般努力就成为义人。如果皇帝的一封信，加上寥寥数语，就能使那些身负无数控告的人获得释放，并使另一些人获得最高荣誉；那么，全能的天主圣神更会将我们从一切邪恶中释放，赐予我们丰厚的正义，使我们充满确信；正如一颗火星落入广阔的大海，会立刻熄灭或消失，被浩渺的水淹没；同样，所有人间的邪恶，一旦落入这神圣泉源的洗池，也比那火星更快、更容易地被淹没，归于无形。有人会说：\Quote{如果洗池除去我们所有的罪，为何不称为赦罪的洗池、或洁净的洗池，而称为重生的洗池？}因为它不仅除去我们的罪，也不仅洁净我们的过犯，而是如同我们重获新生。因为它重新创造和塑造我们，不是再用泥土塑造，而是用另一种元素，即水的本性。因为它不仅洗净器皿，而是完全重新铸造。被洗净的器皿，即使细心擦洗，仍留有原先的痕迹和污秽的残迹；但落入新模具、并藉火焰更新之物，则脱去一切不洁，从熔炉中出来，与新造之物一样闪耀光辉。正如有人取一座因时间、烟尘、灰尘、铁锈而变暗的金像，重新铸造，还给我们一个彻底洁净、闪闪发光的金像；同样，我们这被罪恶铁锈所腐蚀、因过犯积满烟尘、失去天主起初所赐之美的本性，天主把它重新铸造，抛入水中如同抛入模具，并以圣神的恩宠代替火焰，然后使我们焕然一新，光芒四射，足以与太阳争辉；祂粉碎了旧人，塑造了新人，比先前更灿烂。\par

4. 至论这粉碎与这神秘的洁净，古时先知曾隐晦地说：\BibleQuote{你必像打碎陶匠的器皿一样粉碎他们。} 因为前面的话已充分向我们表明这话是关乎信徒的。他说：\BibleQuote{你是我的儿子，我今日生了你。你向我请求，我必将外邦人赐你为基业，将地极赐你为产业。} 你看见他如何提及外邦人的教会，并讲到基督的王国扩展到四方吗？然后他又说：\BibleQuote{你必用铁杖牧养他们，} 不是凶狠的，而是刚强的：\BibleQuote{你必像打碎陶匠的器皿一样粉碎他们。} 且看，这洗濯礼如今更奥秘地被提出来。因为他说的不是陶器，而是陶匠的器皿。但请注意：陶器一旦被打碎，便无法重新塑造，因火使它变得坚硬。但事实上，陶匠的器皿不是陶制的，而是泥制的；因此，即使它们歪斜了，也能藉工匠的技巧轻易地恢复第二次形状。所以，当天主提到一种无法挽救的灾祸时，他说的不是陶匠的器皿，而是陶器；例如，当他愿教导先知和犹太人他将把城市交付给无法挽救的灾祸时，他命先知拿一个陶酒瓶，在众人面前打碎，并说：\BibleQuote{这座城必这样被毁灭，被粉碎。} \BibleRef{耶 19:11}但当他愿向他们表示美好希望时，他把先知带到一座陶器作坊，他不是让他看一个陶器，而是让他看一个在陶匠手中的泥器皿，正倒在地上；并带他到那里说：\BibleQuote{如果这位陶匠拿起并重新塑造他那已跌倒的器皿，难道我岂不更能在你跌倒时恢复你吗？} \BibleRef{耶 18:6}所以，天主不仅能藉重生的洗濯恢复那些由泥土所造的人，也能藉他们的痛悔补赎，使那些已领受圣神之能而又堕落的人回复到原初的状态。但这还不是你们听有关补赎之言的时候，但愿你们永不陷入需要这些补救的境地，而能始终保全你们即将领受的美丽与光辉，毫无玷污。为使你们能永远如此，来，让我们与你们略谈你们的生活方式。因为在摔跤学校里，运动员的跌倒没有危险。因为摔跤是与朋友们进行的，他们在老师的身上操练一切动作。但当比赛的时刻来临，当竞技场敞开，当观众居高临下入座，当裁判到场时，比赛者若变得疏忽，就必然跌倒，并大大蒙羞而退；若他们认真，就赢得冠冕与奖品。同样，在你们身上，这三十天就像某种摔跤学校，既是操练也是练习：让我们从此刻起就学习战胜那邪恶的魔鬼。因为我们就是要脱去衣服与他交手，领洗后就要与他拳击搏斗。让我们从此刻起学习他的抓法，他常从哪一侧攻击，从哪一侧容易威胁我们，这样当比赛来临时，我们才不感到陌生，不会因看到新的摔跤形式而惊慌；而是在我们中间已经练习过，并学会了他的一切手法，就能勇敢地与他进行这些形式的摔跤。因此，他处处习惯威胁我们，但尤其藉着舌头和口。因为对于他的欺骗和毁灭我们，没有比不节制的舌头和不加约束的发言更方便的器官了。由此我们常发生许多过失，由此我们常招致许多严重的控告。舌头的跌倒之易，有人曾说：\BibleQuote{许多人倒在剑下，但不如因舌头跌倒的人多。} \BibleRef{德 28:22}堕落之严重，同一个人又告诉我们说：\BibleQuote{在石板上滑倒比因舌头滑倒还好。} \BibleRef{德 20:18}他所说的意思是：身体跌倒被压碎，还好过说出毁掉灵魂的言语；他不仅说到跌倒，还劝告我们应多加预防，以免失足，这样说：\BibleQuote{为你的口安上门和闩} \BibleRef{德 20:25}，不是叫我们准备门和闩，而是要以极大的安全，把舌头关在可耻的言语之外；又在另一处，在表明我们需要上天的助佑，既与其同行又先于我们自身的努力，以制伏这头野兽后，先知向上主举手说：\Quote{愿我的双手高举如同晚祭；上主，求你在我的口前设置守卫，看守我嘴唇的门户。} 那位先前劝告的人自己又说：\BibleQuote{谁能在我的口前设立守卫，在我的唇上贴上智慧的封条？} \BibleRef{德 22:27}你没有看见吗？每个人都畏惧这些跌倒并为之哀叹，既给予劝告，又祈祷舌头能得到许多警戒。有人说：如果这个器官给我们带来如此大的毁灭，天主为何当初把它安置在我们内呢？因为确实它有大用处，如果我们谨慎，它就只带来用处，而不带来毁灭。请看，那位说过前面话的人说：\BibleQuote{死亡和生命都在舌头的权下。} \BibleRef{箴 18:21}基督也指出同样的事，他说：\BibleQuote{凭你的话，你将被定罪；凭你的话，你将被成义。} \BibleRef{谷 12:27}因为舌头立于中间，随时可供两面使用。你是它的主人。的确，一把剑立于中间，如果你用它攻击你的敌人，这个器官就成为你的救生工具；但如果你将它的打击刺向自己，那么造成你被杀的原因便不是铁的性质，而是你自己的过犯。因此，让我们这样看待舌头：它是立于中间的一把剑；把它磨利，用以控告你自己的罪。不要将打击刺向你的弟兄。为此，天主用双层防御包围它：以牙齿的篱笆和嘴唇的屏障，免得它鲁莽而轻率地发出不当的言语。你说它受不了吗？那就在其内辔制它。藉着牙齿约束它，如同将它的身体交给这些行刑者，让它们咬它。因为宁可在它现在犯罪时被牙齿咬，也不可在将来它因干渴寻一滴水都不得缓解时，无法得到这安慰。实在，它在许多其他方面也常犯罪：嘲笑、亵渎、口出秽语、诽谤、起誓和作伪证。\par

5. 但为了避免我们今天一次说尽所有内容，使你们困惑，我们先提出一条习惯，即关于避免起誓，以此作为序言，坦率地说——如果你们不避免起誓，我指的不是仅仅伪誓，还包括那些因正义而起的誓，我们将不再讨论任何其他主题。因为老师给学生上课时，若未见到前一课牢固记在记忆中，就不会上第二课；而我们若不能清楚表达第一课，就在第一课未完成前灌输其他内容，这实在不合理。这无异于将水倒入有孔的罐中。因此请格外留意，免得你们堵住我们的口。因为这错误是严重的，而且极其严重，因为它看起来并不严重；正因如此我害怕它，因为没有人害怕它。正因如此这病无法治愈，因为它看起来不像病；正如简单的话语不是罪，所以这也不像是罪，而是厚颜无耻地犯下这过犯；若有人质疑，立刻引来嘲笑，不是对因起誓受质疑者的嘲笑，而是对想要纠正这病之人的嘲笑。因此我大量扩展关于这些事的论述。因为我希望拔除深根，消除长久的恶习：我指的不仅是伪誓，还包括出于诚信的起誓。有人说：但某某人，一个忍耐之人，献身于司铎职，生活节制而虔诚，也起誓。不要向我提这个忍耐、节制、虔诚并献身于司铎职的人；如果你愿意，再加上这人是伯多禄或保禄，甚至是从天降下的天使。因为即使如此，我也不顾及他们身份的尊贵。因为我所读关于起誓的法律，不是仆人的法律，而是君王的法律；当宣读君王的谕令时，仆人的一切主张都应静默。但若你能说基督命令我们起誓，或基督没有惩罚起誓的行为，就请指给我看，我便被说服。但若他以如此谨慎禁止此事，并如此关注，以至于将起誓者与恶者并列（因为他说：\Quote{凡是超过这些的，即超过是而非的，便是出于邪恶}）\BibleRef{玛 5:35}，为什么你引述某某人？因为天主不是因你们同仆的疏忽，而是因他自己法律的命令，判定你们有罪。他说：我命令了，你们应当服从，而不是以某某人为庇护，关心他人的恶。既然伟大的达味犯了重罪，我们犯罪就安全吗？请告诉我：那么正因如此我们应当确保这点，只效法圣人的善工；若发现任何疏忽和违反法律之处，我们应当谨慎地逃避。因为我们的算账不是与同仆，而是与我们的主，我们要向他交出一生所行的一切。因此让我们为此审判台做好准备。因为即使违反这法律的人极其尊贵伟大，他必将因过犯受罚。因为天主是不看情面的。那么如何以及以什么方式才能逃避这罪呢？因为不应只指出罪的严重性，还应指明如何逃脱。你有妻子吗？有仆人、儿女、朋友、熟人、邻人吗？要向所有人嘱咐对这些事的谨慎。习惯是可怕的事，难以取代，难以防范，常常在不知不觉中攻击我们；因此，你既然知晓习惯的力量，就应努力摆脱任何坏习惯，转至其他最有益的习惯。因为正如那习惯常能绊倒你，尽管你小心、戒备、思考、考虑；同样，若你转向不宣誓的好习惯，你就无法无意或粗心地落入宣誓的过错中。因为习惯确实强大，具有本性之力。因此，为使我们不不断自寻烦恼，让我们转向另一种习惯，你请求每个亲属和熟人赐予这恩惠：劝告并敦促你逃避宣誓，当被察觉时责备你。因为他们对你的监督，对他们自己也是劝告和向善的提示。因为责备他人宣誓的人，自己不会轻易落入这陷阱。因为大量宣誓不是普通的陷阱，不仅关乎小事，也关乎大事。而我们，无论是买蔬菜，为两文钱争吵，或对仆人大怒并威胁他们，总是呼号天主作证。但一个自由人，拥有某种空虚的尊严，你不敢在市场呼号他为此类事作证；即使你尝试，你也会因你的无礼受罚。然而天上的君王、天使的主，你在为购买和钱财等事争论时，竟把他拉来作证。这如何能容忍呢？那么我们从何处逃脱这恶习？在设下我所说的那些守卫后，让我们为自己定下改过的时间，并加上惩罚，若时间过后仍未改正。多长时间足够呢？我认为非常谨慎、警醒、关心自己救恩的人，不需要超过十天，就能完全摆脱宣誓的恶习。但若十天后我们仍被发现起誓，就给自己加上应得的惩罚，并定下最严厉的处罚和定罪；这定罪是什么呢？我不定，而是让你们自己决定判决。这样我们安排自己的事，不仅关于宣誓，也关于其他缺点，为自己定下时间，并以最严厉的惩罚约束，若任何时候落入其中，就能洁净地来到主前，逃脱地狱之火，并满怀信心地站在基督的审判台前，愿我们所有人都能达成，赖我们的主耶稣基督的恩宠和慈爱，愿光荣归于他及圣父及圣神，世世无穷。阿们。\par

% structure: p0011 source=h2 target=section reason=relative-heading-rank; base=h2

\section{第二篇训诲}

针对那些即将受光照者；以及关于那些以编发、金饰装扮自己的妇女，以及那些曾使用预兆、护符和咒语的人，这一切都与基督信仰相异。\par

1. 我前来首先请求，因着近日出于友爱对你们所说的话语，获得一些果实。因为我们说话，并非只为让你们听听，而是为让你们记住所说的话，并以你们的作为向我们证明这一点。其实，不是向我们，而是向那位洞察人心的天主。正因如此，训诲才得此名称，为使我们即便不在场，我们的言论也能教导你们的心。如果仅在十天的间隔之后，我们就前来要求从所撒的种子中收获果实，请不要惊讶。因为在一天之内，既可以撒种，也可以完成收割。因为我们蒙召战斗，所依靠的不仅是我们自己的力量，更是来自天主的影响。因此，凡领受了所说的话，并藉行为实现了的人，应继续向着前面的事追求。但凡尚未达到这美好成就的人，应立即达到，好能以日后的勤奋消除因怠惰而生的谴责。因为这是可能的，确实有可能，一个极度怠惰的人，若日后勤奋，便能弥补过去时光的全部损失。为此，祂说：\Quote{今天如果你们听从祂的声音，不要心硬，像在激怒之日那样。}祂这样说，是为了劝诫和劝告我们：我们决不可绝望，只要我们还在世上，就应抱有希望，抓住眼前的事，向着从天上召叫我们的奖赏奔跑。那么，就让我们这样做，并探讨这伟大恩赐的名称吧。因为，正如对这尊严的伟大一无所知会使受此尊荣的人更加怠惰，同样，一旦知道，就会使他们感恩，并更加热忱；无论如何，那些从天主那里享受如此荣耀和尊荣的人，竟然不知道这些名称所要表明的是什么，那将是可耻和可笑的。我为何要谈论这恩赐呢？因为如果你们考虑我们族类的普通名称，你们将得到最大的教训和修德的激励。因为对于\Quote{人}这个名称，我们不按照外界之人的定义来界定，而是按照圣经的指示。一个人并非只要有人的手脚，或只拥有理性，而是勇敢地实践虔诚和德行的人。且听祂关于约伯的话。当祂说\Quote{在\Place{乌斯}地方有一个人}，祂并没有像外界之人那样描述他，也不是因为他有两只脚和宽大的指甲，而是加上了他虔诚的证据，说：\Quote{正直、真诚、敬畏天主，远离一切恶事}\BibleRef{约 1:1}，这说明这就是一个人；同样，另一个人说：\Quote{敬畏天主，遵守他的诫命，因为这是人的全部。}\BibleRef{训 12:13}但如果\Quote{人}这个名称提供了如此大的修德激励，那么\Quote{信者}这个名称就更甚了。因为你们被称为信者，是因为你们对天主有信德，并从他那里领受了义德、圣化、灵魂的洁净、义子地位、天国。祂将这一切托付给你们，交给了你们。而你们反过来也将其他事物托付给祂：施舍、祈祷、节制以及一切其他德行。我为何说施舍呢？即使你们给祂一杯凉水，你们绝不会失去它，连这一点祂也会小心地为那一日保存，并丰盛地归还。这确实奇妙，祂不仅保留托付给祂的东西，而且在回报时使之增长。\par

此亦是他按你的能力所命令你，以所托付给你的，延伸你所接受的圣洁，并使从圣洗而来的义德更为光明，使恩宠的恩赐更为灿烂。正如保禄所行，他藉着后来的劳作、热忱和勤勉，增加了他所领受的一切美善。请看天主的谨慎：他当时既没有将全部赐给你，也没有将全部保留，而是赐予一部分，应许了一部分。他为什么不当时就赐予全部呢？为叫你显出对他的信德，单凭他的应许，相信那尚未赐予的。再者，他为什么又不当时全部施予，却赐下了圣神的恩宠、义德和圣化呢？为叫你减轻你的劳苦，并藉着已经赐予的，也使你对那将要来的怀着美好的望德。为此，你即将被称为\Quote{新受光照者}，因为你的光明常新——只要你愿意——永不熄灭。因为白天的光，无论我们愿意与否，黑夜总要接续，但黑暗不认识那光线的照耀。\BibleQuote{光明在黑暗中照耀，黑暗决不能胜过它。}太阳光线出来时，世界尚且不那么明亮，如同灵魂在领受圣神的恩宠并更准确地认识事物的本质时所发出的光明和灿烂。因为当黑夜降临，一片黑暗时，人常看见一团绳索，以为是蛇，因而逃避靠近的朋友如同敌人，又因听见某种响声而大为惊慌。但白昼来到时，就不会发生这类事，一切便显出其本相。在我们灵魂上也是如此。因为恩宠来到，驱散了理智的黑暗，我们便认识了事物的真相，那先前对我们可怕的，如今成为可鄙的。因为藉这神圣的入门礼，我们准确认识了死亡，不再惧怕死亡：知道死亡不是死亡，而是睡眠和适时的安寝。也不怕贫穷、疾病或任何这类事，因为我们知道自己正走向更好的生命，那无玷、不朽、摆脱一切这类变迁的生命。\par

2. 因此，我们不要再贪求今生的事物，也不要贪图餐桌上的奢侈或衣着的华贵。因为你们拥有最卓越的衣裳，你们拥有属神的筵席，你们拥有来自高天的荣耀，基督已成为你们的一切：你们的筵席、你们的衣裳、你们的家园、你们的元首、你们的根干。\BibleQuote{因为凡你们在基督内受洗的，就是穿上了基督。}\BibleRef{迦 3:27}看，祂如何成了你们的衣裳。你们愿意学习祂如何成为你们的筵席吗？\BibleQuote{吃我的人，}祂说，\BibleQuote{就如我因圣父而生活，他也将因我而生活；}以及祂如何成为你们的家园：\BibleQuote{谁吃我的肉，就住在我内，我也住在他内；}\BibleRef{若 6:56}祂又如何是根干，祂又说：\BibleQuote{我是葡萄树，你们是枝条，}\BibleRef{若 15:5}以及祂是弟兄、朋友和新郎：\BibleQuote{我不再称你们为仆人，因为你们是我的朋友；}\BibleRef{若 15:15}\Person{保禄}又说：\BibleQuote{我聘定了你们，把你们如同贞洁的童女献给基督；}\BibleRef{格后 11:2}又说：\BibleQuote{为使祂在众弟兄中作长子；}\BibleRef{罗 8:29}我们不仅成为祂的弟兄，也成为祂的子女：\BibleQuote{看，}祂说，\BibleQuote{我和天主所赐给我的孩子，}\BibleRef{依 8:18}不仅如此，还是祂的肢体和祂的身体。因为如果前面所说的还不足以显明祂向我们显示的慈爱与善意，祂又加上另一件更伟大、更亲近的事，就是称自己为我们的元首。知道了这一切，亲爱的，请以最好的品行回报你们的恩主，并思想这祭献的伟大，装饰你们身体的肢体；思想你们手中所领受的，绝不可让它击打任何人，也不可因击打的罪而使享用如此大恩的事物蒙羞。思想你们手中所领受的，要使之远离一切贪婪和勒索；要想着，你不仅用手领受这圣事，还用口领受，所以要洁净地保守你的舌头，远离下流和无礼的言语、亵渎、伪誓以及所有这类的事。因为，由如此可畏的奥迹所分施、并被如此宝血染红、而成为金剑之物，竟被用于戏弄、侮辱和诙谐，这是何等可悲。要敬畏天主所赐予它的尊荣，不要使它堕落至罪的卑贱；并且再想一想，在手和舌之后，心也领受这可畏的奥迹，所以决不可对邻舍设任何阴谋，而要保守你的思想远离一切邪恶。这样，你也能保守你的眼目和听力安全。因为，在这来自天上的奥秘之声——我是指革鲁宾的呼声——之后，还用淫秽的歌曲和放荡的曲调玷污你的听力，岂不荒唐？如果你用同一双眼目，既观看那不可言喻的、可畏的奥迹，又观看娼妓，并在心中犯奸淫，岂不应受最严厉的惩罚？亲爱的，你被邀请参加婚宴：不要穿着肮脏的衣服进去，而要穿上合宜的婚袍。因为如果人们被邀请参加物质的婚宴时，即使比所有人贫穷，也常常设法获得或购买洁净的衣服，然后去见邀请他们的人。那么，你这被邀请参加属神婚宴和王室盛宴的人，也该思想买怎样的衣服才合适，但甚至无需购买，因为邀请你的人亲自白白地送给你，为使你无法以贫穷为借口。因此，要保存你所领受的衣服。因为如果你丢失了，以后就无法使用，也无法购买。因为这类衣服无处可卖。你们有没有听过，那些古代领受入门圣事的人，如何呻吟捶胸，他们的良心因而激动？所以，亲爱的，要当心，免得你们遭受同样的事。但如果你不抛弃恶人的坏习惯，又怎能不遭受呢？这就是为什么我以前说过，现在还要说，并且不会停止地说：如果谁还没有改正他道德上的缺陷，也没有为自己装备易于获得的德行，就让他不要领受圣洗。因为洗礼池能赦免先前的罪，但有极大的恐惧和危险，以免我们重蹈覆辙，使我们的疗法变成创伤。因为恩宠越大，对于在这些事后犯罪的人的惩罚也越大。\par

3. 因此，为避免我们返回先前的呕吐，让我们从今以后管束自己。因为我们必须先悔改，戒除先前的恶行，然后带着恩宠前来，且要听若翰所说的话，以及宗徒之长对即将受洗者所说的话。因为一位说：\Quote{结出与悔改相称的果实，并不要在心里说：我们有亚巴郎为父。}\BibleRef{路 3:8}；另一位又对询问他的人说：\Quote{你们悔改，并各人以主耶稣基督的名受洗。}\BibleRef{宗 2:38}。凡悔改的人，就不再触碰他曾悔改的同样事情。因此，我们也受命说：\Quote{我弃绝你，撒殚。}，以便我们永不再回到他那里。正如画家写生的情形，愿你们也如此。他们布置画板，描上白线，勾勒出君王的肖像轮廓，在涂上实际色彩之前，会擦去一些线条，修改另一些，自由地纠正错误和不当之处。但当他们最终上色后，就再也不能擦去或更改任何东西，因为会损害肖像的美感，结果成为眼中钉。请思考，你的灵魂就是这幅肖像；因此，在圣神的真色彩来临之前，请擦去那些错误地根植于你心中的习惯，无论是发誓、说谎、傲慢、卑劣的谈吐、戏笑，或是你习惯做的任何非法之事。戒除这习惯，以免你在领洗后又回到它那里。洗礼池使罪恶消失。纠正你的习惯，这样当色彩涂上，君王的肖像显现后，你将来就不会再擦去它们，并给天主赐予你的美加上损伤和疤痕。因此，要抑制愤怒，熄灭情欲。不要烦恼，要同情，不要被激怒，也不要说：\Quote{我的灵魂受了伤害。}没有人能在灵魂上伤害我们，除非我们自己伤害自己的灵魂；这如何可能，我现在就说明。有人夺去了你的财物吗？他没有伤害你的灵魂，而是伤了你的钱财。但你若对他怀恨，你就伤害了自己的灵魂。因为被夺去的财物没有给你造成损害，甚至有益处，但你不消除愤怒，将因这怀恨而在来世交账。有人辱骂你、侮辱你吗？他丝毫未伤害你的灵魂，甚至未伤你的身体。你回骂侮辱了吗？你就伤害了自己的灵魂，因为你要为所说的话在那里交账；我尤其愿你知道这一点：基督徒和忠信的人，无人能伤害其灵魂，连魔鬼自己也不能；不仅天主使我们不受他一切计谋的侵害是奇妙的，而且他还使我们适于修德，并且若我们愿意，即使我们贫穷、身体虚弱、被遗弃、无名、为奴，也毫无阻碍。因为无论贫穷、疾病、身体残缺、奴役，或其他任何这类事，总不能成为修德的阻碍；我为何说贫穷、奴仆、无名呢？即使你是个囚犯，这对修德也绝无阻碍。我接着说这是如何的。你家中有人使你悲伤、激怒你吗？消除你对他的愤怒。镣铐、贫穷、卑微在这件事上阻碍你了吗？我为何说阻碍？它们反而帮助并有助于抑制骄傲。你见他人亨通吗？不要嫉妒他。因为在这种情况下，贫穷也不是障碍。再者，每当你需要祈祷时，以清醒警觉的心去做，在这种情况下，什么也不能阻碍。要显示全部的温良、忍耐、节制、庄重。因为这些事不需要外在的帮助。而这正是修德的要点：它不需要财富、权力、荣耀或任何此类之物，只需一个圣化的灵魂，别无他求。看，同样的事也发生在恩宠上。即使有人瘸腿、被挖去眼睛、身体残缺、陷入极度的衰弱，恩宠也不被任何这些事所阻碍。因为它只寻求一个乐于接受的灵魂，而忽略所有这些外在事物。因为在世俗的军队中，那些将要入伍的人寻求身体的魁梧和健康，而且即将成为士兵的人不仅需要这些，还必须自由。若有人是奴隶，就被拒绝。但天上的君王不寻求任何这类事，却接纳奴隶、老人、肢体虚弱的人进入他的军队，并不以为耻。还有什么比这更慈悲的？什么能更仁慈？因为他寻求我们能力之内的事，而他们寻求我们能力之外的事。因为做奴隶或自由不是我们的行为。高或矮也不是我们能力之内的，年老或强壮等等也是如此。但温良、仁慈等事是我们自己选择的事；天主只要求我们掌管那些事。这是很合理的。因为他召叫我们进入恩宠并非出于自己的需要，而是为施恩于我们；但君王们是因为他们所需的服务；他们把人带入外在的、物质的战争，而他带入属灵的争战；不仅在异教战争中，在竞技中也可看到同样的类比。那些即将被带入场的人，直到传令官亲自在众人面前带他们巡行，高声喊说：\Quote{有人控告这人吗？}尽管那场斗争不涉及灵魂，只涉及身体。那么，你为何要求高贵的证据呢？但在这种情况下，没有这类事，一切都不同，我们的争战不是手扭手，而是灵魂的智慧与心灵的卓越。我们冲突的主席所做的相反。因为他并不带我们巡行说：\Quote{有人控告这人吗？}反而高呼：尽管所有人，尽管魔鬼与撒殚一同站起来，控告他犯有极端和难以言表的罪行，我不拒绝他，也不厌恶他，而是把他从控告者中带走，解除他的邪恶，就这样带他进入竞赛。这是非常合理的。因为在那里，主席对参赛者的胜利毫无贡献，只是站在中间。但在这里，这场圣德竞赛的主席成为战友和助手，与他们一同对抗魔鬼。\par

4. 不仅祂赦免我们的罪是奇妙的，而且祂甚至不揭露它们，也不让它们显露和公开，也不强迫我们走到人群中间，说出我们的过错，而是命令我们只向祂辩护，并向祂告解。然而在世俗的法官那里，如果有人告诉任何一个强盗或盗墓者，当他们被捕时，说出他们的过错便可免于惩罚，他们会非常乐意这样做，因渴望安全而轻视羞耻。但在此事上，并没有这类情况，祂既赦免罪过，也不强迫我们在任何观众面前排列它们。祂只寻求一件事，就是享受这赦免的人应学习这恩赐的伟大。因此，祂在为我们服务时只满足于我们的见证，而在我们为祂服务时却寻求他人作见证，并为了炫耀而行事，这岂不是荒谬吗？当我们惊叹祂的仁慈时，让我们展示我们的行为，并在所有人之前勒住我们舌头的冲动，不要总是说话。\Quote{因为言语多，就难免有过犯。}\BibleRef{箴 10:19} 如果你确实有有用的话要说，就开口。但如果没有必要说话，就保持沉默，因为这样更好。你是个工匠吗？当你坐着工作时，唱圣咏。你不想用嘴唱吗？在心里唱；圣咏是一个伟大的伴侣。这样你不会遭受任何严重的后果，而能坐在你的作坊里如同在隐修院中。因为不是地点的适宜，而是道德的严谨会带给我们宁静。至少保禄在作坊中从事他的手艺时，没有损害自己的德行。\BibleRef{宗 18:3} 所以不要说：我怎么能，身为工匠和穷人，成为哲学家？这正是你可以成为哲学家的原因。因为贫穷远比财富更有利于我们的虔诚，劳作也比闲散更有利；因为财富对不加注意的人甚至是一种障碍。当需要消除愤怒、熄灭嫉妒、克制情欲、祈祷、表现出忍耐和温和、仁慈和爱德时，贫穷何时会成为障碍？因为这些不是靠花钱就能做到的，而是靠展示正确的态度；施舍尤其需要钱，但即使在这方面，贫穷也能更大地闪耀。因为那投入两个小钱的妇人比所有人都穷，却超越了所有人。\BibleRef{路 21:2} 所以我们不要认为财富是什么伟大的东西，也不要认为金子比泥土好。因为物质事物的价值不在于它们的本性，而在于我们对它们的估价。如果有人仔细探究，铁比金子必要得多。因为金子对我们生活的任何需要都没有贡献，而铁却提供了我们大部分的需要，服务于无数技艺；我为什么提到金子和铁的比较？因为这些石头比宝石更必要。因为从宝石中做不出什么有用的东西，而从这些石头中却能建造房屋、城墙和城市。但请你告诉我，从这些珍珠中能得到什么益处？还是说，不会发生什么害处？为了让你戴一粒珍珠坠子，无数穷人被饥饿折磨。你能找到什么借口？什么宽恕？\par

你愿意妆饰你的面容吗？不要用珍珠，而要用端庄和尊贵。这样，你的面容在你丈夫眼中将更充满恩宠。因为另一种妆饰通常会使他陷入猜忌、敌意、争吵和纷争的怀疑，因为没有比一张被怀疑的面孔更令人烦恼的了。但怜悯和端庄的装饰会驱散一切邪恶的怀疑，比任何纽带都更能把你的伴侣吸引到你身边。因为天生的美丽并不能像观看者的心态那样赋予面容如此的美貌，而没有什么比端庄和尊贵更容易产生这种心态；所以，如果一个女人相貌姣好，而她的丈夫对她心怀恶意，她在他的眼中就显得是所有女人中最不值钱的；如果她碰巧不美丽，但丈夫对她心怀善意，她就显得比所有人都更美丽。因为判断不是根据被观看者的本性，而是根据观看者的心态。因此，请用端庄、尊贵、怜悯、仁慈、爱德、对丈夫的深情、宽容、温良、忍耐来妆饰你的面容。这些是德行的色彩。藉着这些，你将吸引天使，而非凡人，成为你的恋人。藉着这些，你将有天主的赞许，当天主接纳你时，他一定会为你赢得你的丈夫。因为人的智慧能使他的面容发光，\BibleRef{训 8:1}，更何况女人的德行能照亮她的面容？如果你认为这是极大的装饰，请告诉我，在那一天珍珠有什么益处？但为什么要谈论那一天呢，因为从当下发生的事就能显明这一切。那么，当那些胆敢辱骂皇帝的人被拖到审判厅，面临极刑危险时，母亲和妻子们摘下项链、金饰、珍珠和一切装饰，脱下华丽衣裳，穿上朴素简陋的衣服，洒上灰烬，匍匐在审判厅门前，就这样赢得了法官的心。如果在这尘世的法庭上，金饰、珍珠和华服会成为陷阱和出卖，而宽容、温良、灰烬、眼泪和粗陋的衣服说服了法官，那么在那公正而可畏的审判台前，岂不更是如此？你将能提出什么理由，什么辩护，当主人用这些珍珠指控你，并把那些因饥饿而死的穷人带到中间时？为此保禄说：\BibleQuote{不以编发、黄金、珍珠或华贵的衣服为装饰。}\BibleRef{弟前 2:9}因为其中会有陷阱。即使我们持续享受它们，我们仍将在死亡时放下它们。但出于德行的一切都是安全的，没有变迁和变化，它在这里使我们更安全，也伴随着我们到那里。你愿意拥有珍珠，永不放下这笔财富吗？摘下一切装饰，通过穷人把它放在基督手中。当他以极大的光辉举起你的身体时，他将为你保存所有的财富。那时，他将赐给你更好的财富和更大的装饰，因为现在的装饰是卑贱而荒谬的。那么，想想你愿意取悦谁，为了谁而戴上这些装饰，不是为了绳匠、铜匠和小贩的赞美。那么，你难道不羞愧，不脸红吗？你为了他们所做的一切，却不认为他们值得你搭话。\par

那么，你将如何嘲笑这种幻想呢？如果你记住你在入门时所说的那句话：\Quote{我弃绝你，撒殚，以及你的浮华和你的侍奉。}因为对珍珠的狂热是撒殚的浮华。你领受金子，不是为了把它绑在身体上，而是为了释放和供养穷人。因此要常常说：\Quote{我弃绝你，撒殚。}没有什么比这句话更安全，如果我们用行为来证明它的话。\par

5. 我认为你们这些即将入门的人应当学习这一点。因为这句话是与主人订立的盟约。正如我们买奴隶时，先问被卖的人是否愿意做我们的仆人，同样，基督在接纳你进入他的服务之前，先问你是否愿意离开那残酷而无情的暴君，并从你那里接受盟约。因为他的服务不是强加给你的。请看天主的仁慈。我们，在付代价之前，先问那些被卖的人，当知道他们愿意时，才付代价。但基督不是这样，他竟为我们所有人付了代价，他的宝血。因为他说：\BibleQuote{你们是被高价买来的。}\BibleRef{格前 7:25}纵然如此，他仍不强迫那些不愿意的人服侍他；他说：除非你蒙恩并自愿，且决心归属我的治下，我不强迫也不勉强。\Emph{我们}本不会选择买邪恶的奴隶。即使有时我们这样选择，也是以扭曲的意愿购买，并付出相应的代价。但基督买那些忘恩负义、无法无天的奴隶时，付出了头等仆人的代价，不，远不止，而且大到言语和思想都无法表达其伟大。因为他既没有给出天堂、大地或海洋，而是给出了比这一切更宝贵的东西，自己的血，这样买了我们。此外，他不要我们提供见证人或登记，只满意于这一句话，只要你从心里说出来。\Quote{我弃绝你，撒殚，以及你的浮华}这句话包含了所有。因此，我们要这样说：\Quote{我弃绝你，撒殚}，如同那些将在那世界的那一天被要求交出这话的人一样，并要保存它，以便那时我们能完好地归还这笔存款。但撒殚的浮华是剧场、马戏场、一切罪恶、守日子、咒语和预兆。\par

有人问：\Quote{什么是征兆？} 常常，当他从自己家里出来时，看见一个独眼人或瘸子，便避开他，当作一个征兆。这是撒殚的排场。因为遇见那人并不会使日子变坏，但生活在罪恶中却会如此。那么，当你出门时，要提防一件事——就是不要让罪恶遇见你。因为正是这罪恶绊倒我们。没有这罪恶，魔鬼不能对我们有任何伤害。你说什么？你看见一个人，便避开他当作征兆，却没有看到魔鬼的陷阱，他如何使你与那没有亏负你的人争战，他如何使你无端与你的弟兄为敌；但天主命令我们爱仇敌；你却避开那没有亏负你的人，对他没有任何指控，难道你不觉得这是多么荒谬，多么可耻，甚至多么危险吗？我能说出更荒谬的事吗？我实在羞愧，脸红；但为了你们的得救，我不得不说出来。如果遇见一个贞女，他就说那天不吉利；但如果遇见一个娼妓，他就说那天吉利、有利、充满许多事务；你羞愧吗？你捶胸顿足，俯伏在地吗？但不要因我所说的言语如此行，而要因所行的事。请看，在这种情况下，魔鬼如何隐藏他的陷阱，好叫我们远离端庄的人，却问候并与不贞洁的人友好。因为他听到基督说：\Quote{凡注视妇女而贪恋她的，心里已与她犯了奸淫。} \BibleRef{玛 5:28} 并且看见许多人胜过了不贞洁，他想要借另一种过犯将他们再次投入罪恶中，便借这迷信的遵守，很高兴地说服他们注意娼妓般的女人。\par

至于那些使用符咒和护身符，并用马其顿的亚历山大的金币围绕他们的头和脚的人，又该说什么呢？告诉我，这些是我们的希望吗？在钉死了我们的主之后，我们竟将一个希腊国王的像当作我们得救的希望？难道你不知道十字架成就了何等伟大的事吗？它消灭了死亡，熄灭了罪恶，使阴府失效，解除了魔鬼的权势，难道它不值得为身体的健康而信靠吗？它复活了整个世界，你还不从中获得勇气吗？告诉我，你配受什么惩罚？你不只随身带着护身符，还把那些醉酒昏聩的老太婆带进屋里，施行咒语，你难道不羞愧、不脸红吗？在如此伟大的哲学之后，竟被这些事情吓倒？还有比这个错误更严重的事。因为当我们提出这些劝告，将他们引开时，他们以为自己在辩护，说那个施行咒语的女人是基督徒，并且除了天主的名字之外什么都不念。正因为此，我特别憎恨并远离她，因为她为了猥亵而使用天主的名字。因为连魔鬼也称呼天主的名字，但它们仍然是魔鬼，并且它们曾这样对基督说：\Quote{我们知道你是谁，你是天主的圣者。} \BibleRef{谷 1:24} 尽管如此，祂还是斥责它们，并把它们赶走。因此，我恳求你们从这错误中洁净自己，并紧握这句话当作手杖；正如你们中间没有人愿意不穿鞋、不披外衣就下到市场去一样，同样，没有这句话，绝不要进入市场，但当你将要跨过门槛时，先说出这句话：我离开你的行列，撒殚，以及你的排场和你的服事，我加入基督的行列。没有这句话，决不要出门。这将是你的手杖，你的盔甲，一座不可攻破的堡垒，并且让这句话伴随着你额上的十字圣号。因为这样，不仅遇见你的人，甚至魔鬼自己，也完全不能伤害你，当他看见你到处带着这些武器出现时；从此以后，你要用这些方法操练自己，好使你领受印记时，成为一个装备精良的士兵，并树立你的战利品以对抗魔鬼，赢得义德的冠冕。愿我们众人都有分获得这冠冕，借着我们的主耶稣基督的恩宠和慈爱，愿光荣归于祂，偕同圣父及圣神，直到永远——阿们。\par

\SourceRecord{Translated by T.P. Brandram.；Nicene and Post-Nicene Fathers, First Series；Vol. 9.；Edited by Philip Schaff.；Buffalo, NY: Christian Literature Publishing Co.,；1889.；<http://www.newadvent.org/fathers/1908.htm>.}
